\section{Research roadmap}
\label{sec:research-roadmap}

The following ordering emphasizes tasks that either validate an existing
claim or unlock several later directions.

\subsection{Upper bounds}

\begin{enumerate}
    \item Carry out the constant optimization described in
    Section~\ref{sec:inverse-sumset-route} using the 2024 high-density
    \(3k-4\) theorem.  Record explicitly whether it improves \(1/3.1955\).
    \item Develop an \(m\)-dependent upper-bound mechanism.  A first target is
    any bound strictly below \(1/3.1955\) for all sufficiently large \(m\); a
    longer-term target is \(1/8+o(1)\).
    \item Strengthen the LP relaxation or study a more informative dual of the
    \(m=3\) interval model.  The naive relaxation has optimum
    \(\min\{N/5,1\}\), so the goal is a reusable valid inequality rather than
    a larger fixed-\(N\) case split.
\end{enumerate}

\subsection{Constructions}

\begin{enumerate}[resume]
    \item Prove or disprove the two-interval formula
    \(M_2(m)=2(m-2)/(m(m+2))\) for \(m\geq4\).
    \item Prove uniformly the two doubled-tail families from
    Section~\ref{sec:extra-intervals}, whose exact endpoint formulas have now
    been checked through \(m=100\).  Then prove or disprove their saturation
    in the fixed-base model, account for the exceptional third interval at
    \(m=12\), and test whether freely moving the base gives a larger gain.
    \item Complete the exact comparison between the dilated finite-group
    construction and the circle construction, including all rounding cases.
\end{enumerate}

\subsection{A possible route beyond fixed \texorpdfstring{$N$}{N}}

Proving \(M_N(3)\leq1/5\) separately for larger \(N\) will not settle the
measurable problem unless the proof is uniform.  A useful theoretical result
would have one of the following forms:
\begin{itemize}
    \item a bounded subconfiguration principle: every union of intervals of
          measure $>1/5$ contains a small subunion whose relevant sumset
          already covers the missing point;
    \item an induction in which bounds for all proper subunions yield a fixed
          collection of inequalities for the full union; or
    \item a compactness/approximation argument paired with a bound independent
          of the number of components.
\end{itemize}
The removable-interaction experiments are relevant precisely when they expose
such a uniform principle.  Merely increasing \(N\) in the MILP is not expected
to do so on its own.
